\section{What the mean-zero assertion says about primes}\label{sec:pnt}
The zero mean also has a two-way elementary connection to the prime
number theorem. The following proof transports an input modulus, so
the equivalence has content independently of Theorem~\ref{thm:strongM}.
Write $\psi(x)=\sum_{p^a\le x,\ a\ge1}\log p$ and
$\Lambda(n)=[n=p^a]\log p$ with the prime recovered by bounded
factorization and value zero if there is no such prime.
Write $m(x)=\sum_{d\le x}\mu(d)/d$ and take $x$ to be a positive
integer unless a real argument is indicated.

\begin{theorem}
The following limits imply one another with PR transformations of any
supplied PR modulus:
\[
  \frac{\psi(x)}x\longrightarrow1,\qquad
  \frac{M(x)}x\longrightarrow0,\qquad m(x)\longrightarrow0.
\]
The first limit is equivalent, in the same sense, to
$\pi(x)\log x/x\longrightarrow1$, where $\pi(x)$ counts primes.
\end{theorem}

\paragraph{From primes to the mean.}
The operation $Df(n)=f(n)\log n$ obeys the product rule for
convolution by $\log(ab)=\log a+\log b$.
Unique factorization gives $D\mathbf1=\mathbf1*\Lambda$.
Apply $D$ to $\mu*\mathbf1=\varepsilon$ and convolve with $\mu$.
It follows that $D\mu=-\mu*\Lambda$. On summing through $x$ we get
\[
 \sum_{n\le x}\mu(n)\log n=-\sum_{d\le x}\mu(d)\psi(x/d).
\]
By the floor identity, $xm(x)=1+\sum_{d\le x}\mu(d)\{x/d\}$,
so $|m(x)|\le1+1/x$. Also
$\sum_{n\le x}\log(x/n)\le\int_0^x\log(x/t)\,dt=x$
by monotonicity (the first term can instead be kept separately).
Given a rational $\eta\in(0,1)$, take $A\ge2$ from the PNT modulus
so that $|\psi(y)-y|\le\eta y$ for $y\ge A$. For real $y$ use
$\psi(y)=\psi(\lfloor y\rfloor)$ and increase the integer cutoff
beyond $2/\eta$, using the initial modulus at $\eta/2$.
For $1\le y<A$ the crude bound is $|\psi(y)-y|\le A^2+A$, by
$\Lambda(n)\le\log n$. Thus
\[
 \frac{|M(x)|}{x}\le
 \frac{A^2+A+3+\eta(1+\log x)}{\log x}\qquad(x\ge2).
\]
Indeed the main term after replacing $\psi(x/d)$ by $x/d$ is at most
$x|m(x)|$, and the remaining large-argument errors cost at most
$\eta x\sum_{d\le x}1/d$. For accuracy $1/k$ take
$\eta=1/(4k)$ and, for example, $x\ge2^{8k(A^2+A+5)}$.

\paragraph{A finite summation-by-parts lemma.}
Suppose $M(x)/x\to0$ with a modulus. For each of the three functions
of a rational $y\ge1$ displayed below one has
\[
 \frac1x\sum_{d\le x}\mu(d)g(x/d)\longrightarrow0:
 \quad
 \begin{cases}
  g_1(y)=\{y\},\\
  g_2(y)=y(H_{\lfloor y\rfloor}-\log y-\gamma),\\
  g_3(y)=\log(\lfloor y\rfloor !)-y\log y+y .
 \end{cases}
\]
Here are bounds that prove the lemma without a general equidistribution
or integration theorem. Each $|g(y)|\le8\sqrt y$. For $g_2$ use
$|H_j-\log j-\gamma|\le1/j$ and
$\log(y/j)\le1/j$ for $j=\lfloor y\rfloor$.
For $g_3$ the increasing function $\log t$ bounds $\log(j!)$ between
$\int_1^j\log t\,dt$ and $\int_1^j\log t\,dt+\log j$.
Replacing the upper integration endpoint $j$ by $y$ costs at most
$\log y$. Thus $|g_3(y)|\le 1+2\log y\le8\sqrt y$.
On $[1,A]$, $A$ integral, the two endpoint magnitudes plus the total
variation of each $g$ are bounded by $V_A=50(A+1)^4$.
For a direct check, on the open unit intervals the three derivatives are
$1$, $H_j-\log y-\gamma-1$, and $-\log y$. Their absolute values are
at most $3(A+1)$. The jumps at integers have magnitudes at most
$1,1,\log A$, respectively. The same bounds hold for finite differences
by integrating the explicit derivatives on each intervening piece.

The terms $d\le x/A$ in the original sum contribute at most
$16x/\sqrt A$, using $\sum_{d\le u}d^{-1/2}\le2\sqrt u$.
For the other terms, partial summation with the bounded variation on
$[1,A]$ costs at most $\eta x V_A$ if $|M(t)|\le\eta t$ for all
integers $t\ge\lfloor x/A\rfloor$. Thus for a target $\epsilon>0$
choose $A\ge(32/\epsilon)^2$, then $\eta=\epsilon/(2V_A)$, and take
$x\ge A(U+1)$ where $U$ is the input modulus at $\eta$.
All variation estimates only need finite lists of rational input points
and the finitely many integer breakpoints.

\paragraph{From the mean back to primes.}
The floor identity and the lemma for $g_1$ give $m(x)\to0$.
Define $L_x=\sum_{d\le x}\mu(d)\log(x/d)/d$. Grouping by $dk$ gives
\[
 1=\sum_{d\le x}\frac{\mu(d)}d H_{\lfloor x/d\rfloor}
  =L_x+\gamma m(x)+x^{-1}\sum_{d\le x}\mu(d)g_2(x/d).
\]
Hence $L_x\to1$. Finally $\log n=\sum_{r\mid n}\Lambda(r)$, and
inversion, then summation over $n\le x$, gives
$\psi(x)=\sum_{d\le x}\mu(d)\log(\lfloor x/d\rfloor!)$.
Using $g_3$ gives
$\psi(x)/x=L_x-m(x)+x^{-1}\sum_{d\le x}\mu(d)g_3(x/d)\to1$.
A modulus for $m(x)\to0$ alone gives $M(x)/x\to0$ by
\[
 M(x)=xm(x)-\sum_{j=1}^{x-1}m(j).
\]
Indeed if $|m(j)|\le\eta$ for $j\ge U$, dividing by $x$ costs at
most $2\eta+2U/x$, using the global bound $|m(j)|\le2$.

\paragraph{The prime-counting formulation.}
Put $\vartheta(x)=\sum_{p\le x}\log p$. Prime powers of exponent
at least two contribute at most $\sqrt x\log x$ to $\psi(x)$:
each prime $p\le\sqrt x$ contributes at most $\log x$.
For $r\ge2$ put $y=x^{1-1/r}$ (the cutoff can be compared by integer
$r$th powers). Splitting primes at $y$ gives the elementary bounds
\[
 (1-1/r)\big(\pi(x)\log x-y\log x\big)
       \le\vartheta(x)\le \pi(x)\log x .
\]
The error $y\log x/x$ is at most $2r x^{-1/(2r)}$ since
$\log u\le u$ for $u\ge1$. Choosing $1/r$ below a quarter of the desired
error and $x$ above a suitable integer power of $2r/\epsilon$
gives each implication between the normalized $\vartheta$ and
$\pi$ limits. Taking a finer fixed fraction of the target precision
absorbs the factor $(1-1/r)^{-1}\le2$. These are PR modulus
transformations. The same power bound controls $\log x/\sqrt x$.
